%% kapa: exam=1 round=19 year=2008 period=2 form=A subject=english subject_ko=영어 seq=9 number=49
%% source: Q-Net 감정평가사 기출문제 (cst00309) artlSeq=1001962 file=extracted/1008/1001962/2008년 제19회 감정평가사 1차시험 기출문제/2교시.hwp
%% conversion: kordoc 4.24.51 -> tools/gen1.py (auto); review=proofread-sonnet; flags=none
\begin{question}{49}
% [공통 지문/지시문: 원문에서 여러 문항에 걸쳐 제시됨]
\begin{framedtext}
※ 빈칸에 들어갈 가장 알맞은 것을 고르시오. (47-50) \par
\end{framedtext}
Chinese characters were used for so long as a proper method of expressing Koreans' thoughts and feelings that they \blank{} the intellectual and cultural lives of the Korean people even after the appearance of Hangeul.
\begin{choices}
  \item resulted
  \item affected
  \item inhibited
  \item wielded
  \item promulgated
\end{choices}
\end{question}
\answer{2}
